\chapter{Analyse et conception}
\label{chap:conception}

\callout{En bref}{Ce chapitre formalise les besoins, présente l'architecture
en couches du système, détaille le serveur MCP et son registre d'outils
permissionné, décrit les sept agents et leur matrice de permissions, expose le
moteur déterministe adaptatif, et illustre la conception par des diagrammes UML.}

\section{Spécification des besoins}

\subsection{Besoins fonctionnels}
Le système doit permettre à un utilisateur de :
\begin{enumerate}
  \item \textbf{téléverser} un jeu de données tabulaire (CSV, Excel, JSON) ;
  \item déclencher une \textbf{analyse automatique} de bout en bout ;
  \item suivre en \textbf{temps réel} l'avancement du traitement par étape ;
  \item consulter un \textbf{tableau de bord interactif} (indicateurs et
        graphiques) ;
  \item consulter et \textbf{exporter} un \textbf{rapport décisionnel}
        (page web, PDF, HTML) ;
  \item retrouver et rouvrir les \textbf{exécutions passées} et leurs
        artefacts.
\end{enumerate}

\subsection{Besoins non fonctionnels}
\begin{description}
  \item[Généralité.] Le système ne doit faire aucune hypothèse sur le domaine
    métier ni sur les noms de colonnes.
  \item[Fiabilité.] Une exécution doit toujours produire des livrables
    complets, même en cas de défaillance d'un agent.
  \item[Véracité.] Les chiffres du rapport doivent être vérifiables face aux
    indicateurs effectivement calculés.
  \item[Sûreté.] Chaque agent ne doit accéder qu'aux outils nécessaires à son
    rôle (moindre privilège).
  \item[Reproductibilité.] Chaque exécution doit être tracée dans un magasin
    d'artefacts autonome.
  \item[Observabilité.] La consommation (jetons, latence, coût estimé) doit
    être mesurée et restituée.
  \item[Portabilité.] Le système doit fonctionner avec un fournisseur de LLM
    hébergé (Groq) ou local (Ollama).
\end{description}

\subsection{Acteurs et cas d'usage}
L'acteur principal est l'\textbf{Analyste / Décideur}, qui interagit avec la
plateforme via une interface web. La figure~\ref{fig:usecase} présente le
diagramme de cas d'usage.

\begin{figure}[ht]
\centering
\begin{tikzpicture}[
  uc/.style={ellipse,draw=accent,fill=accentsoft,text=ink,minimum width=4.3cm,
    minimum height=1.05cm,align=center,font=\small,inner sep=1pt},
  line/.style={accentdark,thick}]
  % --- Acteur (stick figure soigné) ---
  \begin{scope}[shift={(-4.4,-3.0)},accentdark,line width=1.2pt]
    \draw[fill=accentsoft] (0,1) circle (0.34);        % tête
    \draw (0,0.66) -- (0,-0.55);                        % tronc
    \draw (-0.55,0.28) -- (0.55,0.28);                  % bras
    \draw (0,-0.55) -- (-0.45,-1.3);                    % jambe gauche
    \draw (0,-0.55) -- (0.45,-1.3);                     % jambe droite
  \end{scope}
  \node[font=\small\bfseries,text=ink] at (-4.4,-4.55) {Analyste};
  \node[font=\scriptsize,text=muted] at (-4.4,-4.95) {/ Décideur};

  % --- Frontière du système ---
  \node[draw=accent,line width=1pt,rounded corners=8pt,fill=accent!3,
        minimum width=7.8cm,minimum height=8.4cm,
        label={[font=\small\bfseries,text=accentdark,yshift=-2pt]north:Plateforme BI multi-agents}]
        (sys) at (2.2,-3.0) {};

  % --- Cas d'usage (répartis verticalement) ---
  \node[uc] (u1) at (2.2,-0.5) {Téléverser un jeu de données};
  \node[uc] (u2) at (2.2,-1.75) {Lancer une analyse automatique};
  \node[uc] (u3) at (2.2,-3.0) {Suivre la progression};
  \node[uc] (u4) at (2.2,-4.25) {Consulter le tableau de bord};
  \node[uc] (u5) at (2.2,-5.5) {Consulter / exporter le rapport};

  % --- Liens acteur -> cas d'usage ---
  \foreach \u in {u1,u2,u3,u4,u5} \draw[line] (-3.7,-3.0) -- (\u.west);
\end{tikzpicture}
\caption{Diagramme de cas d'usage de la plateforme.}
\label{fig:usecase}
\end{figure}

\section{Architecture globale}
L'architecture adopte une organisation \textbf{en couches} (figure~\ref{fig:archi}),
qui sépare nettement les responsabilités et permet une évolution
indépendante de chaque niveau.

\begin{figure}[ht]
\centering
\begin{tikzpicture}[
  L/.style={rounded corners=5pt, draw=accent, line width=0.8pt,
    minimum width=10.2cm, minimum height=1.05cm, text width=9.3cm,
    align=left, font=\small, inner xsep=11pt, inner ysep=6pt},
  down/.style={-{Stealth[length=3mm]}, accentdark, line width=1.4pt},
  back/.style={-{Stealth[length=3mm]}, accent, line width=1pt, dashed}]

  \node[L, fill=accent, text=white] (ui)
    {\textbf{Couche présentation}\\[1pt]{\scriptsize Interface web FastAPI — téléversement, suivi, tableau de bord, rapport}};
  \node[L, fill=accentsoft, below=7mm of ui] (orch)
    {\textbf{Couche orchestration}\\[1pt]{\scriptsize Pipeline + filets de sécurité déterministes}};
  \node[L, fill=accentsoft, below=7mm of orch] (ag)
    {\textbf{Couche agents}\\[1pt]{\scriptsize Orchestrateur · Ingestion · Data Prep · KPI · Dashboard · Publication · Reporter}};
  \node[L, fill=accentdark, text=white, below=7mm of ag] (mcp)
    {\textbf{Serveur MCP}\\[1pt]{\scriptsize Registre de 14 outils + contrôle d'accès par rôle (refus par défaut)}};
  \node[L, fill=accentsoft, below=7mm of mcp] (core)
    {\textbf{Moteur déterministe}\\[1pt]{\scriptsize Profilage · KPIs · Graphiques · Insights · Validation}};
  \node[L, fill=white, below=7mm of core] (store)
    {\textbf{Magasin d'artefacts}\\[1pt]{\scriptsize \texttt{dashboard.html} · \texttt{report.html} · \texttt{state.json}}};

  % Flux descendant : la requête traverse les couches
  \foreach \a/\b in {ui/orch, orch/ag, ag/mcp, mcp/core, core/store}
    \draw[down] (\a) -- (\b);

  % Flux remontant : les résultats reviennent vers l'interface
  \draw[back] (store.east) .. controls +(2.2,0) and +(2.2,0) .. (ui.east)
    node[pos=0.5, fill=white, inner sep=1.5pt, text=muted, font=\scriptsize] {résultats};
\end{tikzpicture}
\caption{Architecture en couches du système.}
\label{fig:archi}
\end{figure}

Le principe directeur est une \textbf{séparation des préoccupations} :
\begin{itemize}
  \item la \textbf{présentation} (interface web) ne contient aucune logique
        d'analyse ;
  \item l'\textbf{orchestration} séquence les agents et garantit la complétude
        des résultats ;
  \item les \textbf{agents} portent le jugement métier (curation, narration) ;
  \item le \textbf{serveur MCP} médie tous les accès aux capacités et applique
        les permissions ;
  \item le \textbf{moteur déterministe} effectue tous les calculs ;
  \item le \textbf{magasin d'artefacts} assure la traçabilité.
\end{itemize}

\section{Le serveur MCP et le registre d'outils}
Le serveur MCP est le \textbf{point de passage obligé} entre les agents et les
capacités du système. Il expose 14 \textbf{outils}, chacun décrit par un nom,
une description en langage naturel et un \textbf{schéma JSON} de ses arguments
— description directement exploitable par les LLM pour décider de leurs appels.
Le tableau~\ref{tab:outils} récapitule ces outils.

\begin{table}[ht]
\centering\small
\caption{Les outils exposés par le serveur MCP.}
\label{tab:outils}
\begin{tabularx}{\linewidth}{@{}lY@{}}
\toprule
\textbf{Outil} & \textbf{Rôle} \\
\midrule
\texttt{load\_dataset} & Charge un fichier (CSV/Excel/JSON) et l'enregistre. \\
\texttt{validate\_dataset} & Contrôle structurel (doublons, colonnes vides…). \\
\texttt{profile\_dataset} & Profile les colonnes (rôles, cardinalité, statistiques). \\
\texttt{clean\_dataset} & Applique des opérations de nettoyage. \\
\texttt{suggest\_kpis} / \texttt{compute\_kpis} & Propose puis calcule les indicateurs. \\
\texttt{detect\_anomalies} & Détecte les valeurs aberrantes et pics temporels. \\
\texttt{compare\_segments} & Compare une mesure entre groupes. \\
\texttt{suggest\_charts} / \texttt{build\_chart} & Propose puis construit les graphiques. \\
\texttt{build\_dashboard} & Assemble le tableau de bord interactif. \\
\texttt{publish\_artifacts} & Fige l'état de l'exécution sur disque. \\
\texttt{generate\_report} & Produit le rapport décisionnel (HTML). \\
\texttt{get\_run\_state} / \texttt{set\_run\_state} & Lit / écrit l'état partagé de l'exécution. \\
\bottomrule
\end{tabularx}
\end{table}

\section{Les sept agents et la matrice de permissions}
Le système répartit le travail entre \textbf{sept agents spécialisés}. Chaque
agent possède un \textbf{rôle}, une \textbf{invite système} (\textit{prompt})
qui le cadre, et un \textbf{ensemble d'outils autorisés} défini dans le
fichier de configuration \texttt{config.yaml}. Le serveur MCP applique ces
permissions à chaque appel, selon un principe de \textbf{refus par défaut}.
Le tableau~\ref{tab:permissions} présente cette matrice.

\begin{table}[ht]
\centering\small
\caption{Matrice rôles / permissions (extraite de \texttt{config.yaml}).}
\label{tab:permissions}
\begin{tabularx}{\linewidth}{@{}lY@{}}
\toprule
\textbf{Agent} & \textbf{Outils autorisés} \\
\midrule
\textbf{Orchestrateur} & \texttt{list\_agents}, \texttt{get\_run\_state}, \texttt{set\_run\_state} \\
\textbf{Ingestion} & \texttt{load\_dataset}, \texttt{validate\_dataset} \\
\textbf{Data Prep} & \texttt{profile\_dataset}, \texttt{clean\_dataset} \\
\textbf{KPI} & \texttt{profile\_dataset}, \texttt{suggest\_kpis}, \texttt{compute\_kpis}, \texttt{detect\_anomalies}, \texttt{compare\_segments} \\
\textbf{Dashboard} & \texttt{profile\_dataset}, \texttt{suggest\_charts}, \texttt{build\_chart} \\
\textbf{Publication} & \texttt{build\_dashboard}, \texttt{publish\_artifacts} \\
\textbf{Reporter} & \texttt{get\_run\_state}, \texttt{generate\_report} \\
\bottomrule
\end{tabularx}
\end{table}

\subsection{Rôles des agents}
\begin{description}
  \item[Orchestrateur.] Planifie l'exécution, observe la forme des données et
    signale les risques. Il ne manipule pas directement les données.
  \item[Ingestion.] Charge le fichier brut et valide sa structure.
  \item[Data Prep.] Profile les colonnes et applique un nettoyage minimal et
    sûr.
  \item[KPI.] Sélectionne, parmi des candidats adaptatifs, les indicateurs les
    plus pertinents, les calcule, puis détecte anomalies et écarts de segments.
  \item[Dashboard.] Choisit un ensemble équilibré de graphiques et les
    construit.
  \item[Publication.] Assemble le tableau de bord HTML et fige les artefacts.
  \item[Reporter.] Rédige le rapport décisionnel en s'appuyant exclusivement
    sur l'état réel de l'exécution.
\end{description}

\subsection{Boucle d'exécution d'un agent}
Chaque agent fonctionne selon une \textbf{boucle d'appel d'outils} bornée :
le LLM reçoit l'invite et la liste de ses outils autorisés, décide d'un ou
plusieurs appels, le serveur MCP exécute et renvoie les résultats, puis le
modèle poursuit jusqu'à produire une réponse finale. Un mécanisme de
\textbf{relance avec incitation} (\textit{retry-with-nudge}) corrige le cas
fréquent où le modèle répond par du texte sans appeler ses outils.

\section{Le moteur déterministe adaptatif}
Le cœur « intelligent » généraliste ne peut reposer entièrement sur le LLM :
face à des dizaines de colonnes, le modèle hallucinerait des noms et
manquerait des agrégations évidentes. Le \textbf{moteur déterministe} produit
donc une base adaptative riche, que les agents \emph{curent} et \emph{narrent}.

\subsection{Profilage par rôles}
Le profileur classe chaque colonne dans un \textbf{rôle sémantique} :
\textbf{mesure} (numérique agrégeable), \textbf{dimension} (catégorielle de
faible cardinalité), \textbf{date}, \textbf{identifiant}, \textbf{texte},
\textbf{booléen}, ou \textbf{géographie} (valeurs reconnues comme des pays).
Ces rôles, déduits par des heuristiques paramétrables, permettent à tout le
reste du système d'agir de manière générique.

\subsection{Génération adaptative}
À partir des seuls rôles, le moteur dérive :
\begin{itemize}
  \item des \textbf{KPIs} candidats (totaux, moyennes, comptages distincts,
        variation période sur période, part du segment dominant) ;
  \item des \textbf{graphiques} candidats (séries temporelles, barres top-N,
        distributions, nuages de points, parts, carte choroplèthe, heatmap de
        corrélations) ;
  \item des \textbf{insights} (anomalies, comparaisons de segments,
        corrélations).
\end{itemize}

\section{Diagrammes de conception}
La figure~\ref{fig:sequence} présente le diagramme de séquence simplifié d'une
exécution, mettant en évidence l'orchestration des sept agents via le serveur
MCP.

\begin{figure}[ht]
\centering
\begin{tikzpicture}[font=\scriptsize,
  life/.style={thick,muted,dashed},
  msg/.style={-{Stealth[length=2mm]},accentdark},
  head/.style={rectangle,rounded corners=3pt,draw=accent,fill=accentsoft,
    text=ink,minimum width=1.7cm,minimum height=0.7cm,align=center,font=\scriptsize\bfseries}]
  % lifelines
  \node[head] (U) at (0,0) {Utilisateur};
  \node[head] (P) at (2.6,0) {Pipeline};
  \node[head] (A) at (5.2,0) {Agents};
  \node[head] (M) at (7.8,0) {Serveur MCP};
  \node[head] (C) at (10.4,0) {Moteur};
  \foreach \n in {U,P,A,M,C} \draw[life] (\n.south) -- ++(0,-7.4);
  % messages
  \def\y{-0.7}
  \draw[msg] (U.south|-0,\y) -- node[above]{téléverse} (P.south|-0,\y);
  \def\y{-1.5}\draw[msg] (P.south|-0,\y) -- node[above]{lance étape} (A.south|-0,\y);
  \def\y{-2.3}\draw[msg] (A.south|-0,\y) -- node[above]{appelle outil} (M.south|-0,\y);
  \def\y{-3.1}\draw[msg] (M.south|-0,\y) -- node[above]{exécute} (C.south|-0,\y);
  \def\y{-3.9}\draw[msg] (C.south|-0,\y) -- node[above]{résultat} (M.south|-0,\y);
  \def\y{-4.7}\draw[msg] (M.south|-0,\y) -- node[above]{renvoie} (A.south|-0,\y);
  \def\y{-5.5}\draw[msg] (A.south|-0,\y) -- node[above]{état mis à jour} (P.south|-0,\y);
  \def\y{-6.6}\draw[msg] (P.south|-0,\y) -- node[above]{artefacts (dashboard, rapport)} (U.south|-0,\y);
  \node[text=muted] at (5.2,-7.0) {\footnotesize\itshape Répété pour les 7 agents : Ingestion → Orchestrateur → Data Prep → KPI → Dashboard → Publication → Reporter};
\end{tikzpicture}
\caption{Diagramme de séquence simplifié d'une exécution du pipeline.}
\label{fig:sequence}
\end{figure}

\subsection{Vue en composants}
La figure~\ref{fig:composants} présente l'organisation en composants
logiciels et leurs dépendances. Le pipeline dépend des agents, qui dépendent
tous du serveur MCP ; celui-ci est le seul à dépendre du moteur déterministe
et du client LLM. Cette dépendance \emph{en entonnoir} matérialise le rôle de
médiateur unique du serveur MCP.

\begin{figure}[ht]
\centering
\begin{tikzpicture}[
  comp/.style={rectangle,draw=accent,fill=white,text=ink,rounded corners=3pt,
    minimum height=0.95cm,minimum width=2.7cm,align=center,font=\scriptsize\bfseries},
  pkg/.style={rectangle,draw=accentdark,fill=accentsoft,text=ink,rounded corners=3pt,
    minimum height=0.95cm,minimum width=2.7cm,align=center,font=\scriptsize\bfseries},
  dep/.style={-{Stealth[length=2.2mm]},muted,thick}]
  \node[comp] (ui) at (0,4) {ui\\(FastAPI)};
  \node[comp] (pipe) at (0,2.6) {pipeline.py};
  \node[comp] (agents) at (0,1.2) {agents\\(7 agents)};
  \node[pkg] (mcp) at (0,-0.3) {mcp\_server\\(outils + perms)};
  \node[pkg] (core) at (-2.9,-2.0) {core\\(moteur)};
  \node[pkg] (llm) at (2.9,-2.0) {llm\\(Groq/Ollama)};
  \node[comp] (store) at (0,-3.6) {Magasin\\d'artefacts};
  \draw[dep] (ui)--(pipe); \draw[dep] (pipe)--(agents); \draw[dep] (agents)--(mcp);
  \draw[dep] (mcp)--(core); \draw[dep] (mcp)--(llm);
  \draw[dep] (agents.east) to[bend left=30] (llm.north);
  \draw[dep] (core)--(store); \draw[dep] (mcp)--(store);
\end{tikzpicture}
\caption{Diagramme de composants et dépendances.}
\label{fig:composants}
\end{figure}

\callout{Synthèse du chapitre}{L'architecture en couches isole présentation,
orchestration, agents, serveur MCP, moteur et artefacts. Le serveur MCP
centralise 14 outils et applique une matrice de permissions stricte aux sept
agents. Le moteur déterministe garantit la généralité ; les agents apportent le
jugement. Le chapitre suivant détaille la réalisation.}
